\chapter{Liver Transplantation}
\section*{Overview}
Liver transplantation replaces a diseased liver with a deceased-donor liver or a portion of a living donor's liver.
\section*{Common Instruments and Equipment}
The exact setup varies by technique and setting. Common items may include:
\begin{itemize}[nosep,leftmargin=*]
  \item transplant surgery instrument set and bypass equipment as needed
  \item vascular clamps and anastomosis sutures
  \item organ preservation solution and temperature-controlled transport container
\end{itemize}

\section*{Why It Is Done}
It is used for selected patients with irreversible end-stage liver disease, acute liver failure, or certain liver cancers.
\section*{How It Is Performed}
The diseased liver is removed and the donor liver is connected to the recipient's major blood vessels and bile duct. Complex vascular and biliary reconstruction is required.
\section*{Contraindications}
No single contraindication applies to every patient; the treating team assesses whether the procedure is suitable for the individual and the reason for treatment.

\section*{Precautions}
The care team reviews relevant health conditions, medicines, allergies, and procedure-specific risks, and explains any changes needed for the individual.

\section*{Risks and Follow-up}
Bleeding, thrombosis, bile leak, rejection, infection, kidney injury, graft dysfunction, and cancer recurrence are possible. Lifelong immunosuppression and specialist monitoring are required.
\section*{References}
\begin{enumerate}\item MedlinePlus. Liver Transplant. U.S. National Library of Medicine; 2025.\item Current Medical Diagnosis and Treatment. McGraw-Hill; 2025.\end{enumerate}
\clearpage
